\section{Estudio de Vigas}

Se presenta el prediseño de las vigas del edificio en función de la luz de 
cálculo y las restricciones arquitectónicas de cada nivel.

\subsection{Criterios de Prediseño}

La altura mínima de las vigas se estima a partir de la relación luz y el alto de la sección:

\begin{equation}
    h \geq \frac{L}{n}
    \label{eq:viga_peralte}
\end{equation}

donde $L$ corresponde a la luz libre de la viga y $n$ es un coeficiente que 
depende de las condiciones de apoyo: se utiliza $n = 10$ para vigas simplemente 
apoyadas y $n = 15$ para vigas empotradas o continuas. El ancho mínimo adoptado en estacionamientos es de 25~cm, 
condicionado por el espesor de los muros de apoyo.


Las vigas de fachada son elementos invertidos cuya altura queda 
definida por los planos de arquitectura, los cuales fijan las cotas de 
cierre de fachada en cada tramo. Se distinguen tres alturas según el sector:
74~cm y 164~cm. Además, se incorpora una viga invertida de 72~cm en el sector de la caja escala.

\subsection{Prediseño por Nivel}

\subsubsection{Subterráneo}

Las vigas del subterráneo son normales, con una restricción de altura de 50~cm. La luz crítica medida en planta es de 8{,}15~m; aplicando 
el criterio para vigas continuas ($n = 15$) en la ecuación~(\ref{eq:viga_peralte}):

\[
    h \geq \frac{8{,}15}{15} = 0{,}40 \text{ m} = 54.33 \text{ cm}
\]

Se adopta una viga de \textbf{25/50~cm} (ancho/alto). Ya que la altura máxima es de 50~cm.

\subsubsection{Primer Piso}

Las vigas de fachada son invertidas con altura definida por arquitectura:

\begin{itemize}
    \item \textbf{Viga 25/74~cm}: viga invertida de fachada.
    \item \textbf{Viga 25/164~cm}: viga invertida de fachada.
    \item \textbf{Viga 25/7~cm}: viga invertida de la caja escala.
\end{itemize}


\subsubsection{Segundo Piso}

Mantiene la misma configuración de fachada que el primer piso. 

\begin{itemize}
    \item \textbf{Viga 25/74~cm}: viga invertida de fachada.
    \item \textbf{Viga 25/164~cm}: viga invertida de fachada.
    \item \textbf{Viga 25/7~cm}: viga invertida de la caja escala.
\end{itemize}


\subsubsection{Piso Tipo (Pisos 3 al 21)}

La planta tipo replica la configuración del segundo piso en los 19 niveles 
restantes: vigas de fachada 25/74~cm y 25/164~cm, viga de escalera 25/72~cm.

\subsection{Resumen de Vigas Tipo por Nivel}

Los resultados del prediseño se resumen en la Tabla~\ref{tab:vigas_tipo}. 
Las dimensiones se expresan como ancho/alto total para las vigas 
invertidas de fachada el peralte indicado corresponde a la altura total 
del elemento.

\begin{table}[H]
    \centering
    \begin{tabular}{llc}
        \toprule
        \textbf{Nivel} & \textbf{Tipo de Viga} & \textbf{Ancho/Alto [cm]} \\
        \midrule
        Subterráneo 
            & Viga subterráneo & 25/50 \\
        \midrule
        \multirow{2}{*}{Primer Piso}
            & Viga invertida de fachada  & 25/74  \\
            & Viga invertida de fachada       & 25/164 \\
            & Viga invertida sector escalera & 25/72  \\
        \midrule
        \multirow{3}{*}{Segundo Piso}
            & Viga invertida de fachada  & 25/74  \\
            & Viga invertida de fachada      & 25/164 \\
            & Viga invertida sector escalera             & 25/72  \\
        \midrule
        \multirow{3}{*}{Piso Tipo (3--21)}
            & Viga invertida de fachada  & 25/74  \\
            & Viga invertida de fachada      & 25/164 \\
            & Viga invertida sector escalera             & 25/72  \\
        \bottomrule
    \end{tabular}
    \caption{Resumen de vigas tipo por nivel}
    \label{tab:vigas_tipo}
\end{table}
